\documentclass[11pt]{article}
\usepackage[margin=1in]{geometry}
\usepackage{amsmath,amssymb,amsthm,booktabs}
\usepackage[colorlinks=true,linkcolor=blue,citecolor=blue,urlcolor=blue]{hyperref}
\hypersetup{pdftitle={A finite-seed parameter refinement for one-sample matroid prophet inequalities},pdfauthor={}}
\newtheorem{theorem}{Proposed theorem}
\newtheorem{lemma}{Lemma}
\newtheorem{proposition}{Proposition}
\DeclareMathOperator{\cl}{cl}
\DeclareMathOperator{\rk}{r}
\DeclareMathOperator{\OPT}{OPT}
\newcommand{\E}{\mathbb E}
\newcommand{\Prb}{\mathbb P}
\newcommand{\one}{\mathbf 1}
\newcommand{\Src}{\cite{source}}
\allowdisplaybreaks
\title{A finite-seed parameter refinement for\protect\\one-sample matroid prophet inequalities}
\author{}
\date{Preparatory research note, October 10, 2026}
\begin{document}
\maketitle
\begin{abstract}
We give a candidate refinement of the constants in OpenAI's one-sample
matroid prophet construction. Retaining its combinatorial machinery,
we derive a coefficient greater than $2^{-63}$ against an adversary
that sees every sample, online value, and seed bit. The accessible source
already derives $2^{-293}$, stronger than its advertised $2^{-310}$.
The changes are a sharper use of existing exit and residual estimates,
affine mixing of branches, and sixteen shared rounding offsets. An
explicit implementation uses $62n+40$ fair bits; no polynomial running
time is asserted. Exact scalar checks and finite online-rule checks
accompany the derivation. This is an AI-assisted, unreviewed draft,
not a formalized proof. Daniel's October 8 announcement of $2^{-47}$
is stronger; its proof was unavailable in our audit. We make no claim
to the strongest current result or to priority.
\end{abstract}

\section{Model, claim, and provenance}

Let $M=(E,\mathcal I)$ be a known finite labeled matroid, $n=|E|$,
with a fixed total label order. For $A,P\subseteq E$, write
$\rk(A\mid P)=\rk(A\cup P)-\rk(P)$ and $L=\cl(\varnothing)$.
For a nonnegative vector $w$, put
$\OPT_M(w)=\max_{I\in\mathcal I}\sum_{e\in I}w_e$.

\begin{theorem}\label{thm:main}
For every such $M$ there is a measurable one-sample online rule whose
seed consists of $62n+40$ independent fair bits and whose law does not
depend on the value distributions. Let all $2n$ variables $S_e,V_e$
be mutually independent, nonnegative, with $S_e$ and $V_e$ identically
distributed for each $e$. Let the seed $R$ be independent of them and
assume $\E\OPT_M(V)<\infty$. For every measurable arrival permutation
$\pi(S,V,R)$, the rule accepts an independent set at every prefix and
its reward satisfies
\[
 \E\operatorname{reward}\ \ge\ 2^{-63}\E\OPT_M(V).
\]
The rule receives all samples initially, uses no distribution description,
and observes online values only upon their arrival. The adversary may
observe the entire initial seed, including unused coins.
\end{theorem}

The argument below depends on named combinatorial lemmas of the
September 23 manuscript \Src, pinned at commit
\texttt{fd4aeeb2ee4fc729c18d98444fed42fd0529eeeb}.
We restate the interfaces needed for retuning and derive all changed
quantitative estimates. We do not claim a new density decomposition,
marked-pivot theorem, reverse-exposure method, or sample/value reduction.
Source arguments, rather than their hard-coded numerical corollaries,
are the dependencies. The source's final calculation is
$2^{-140-16-5-100-32}=2^{-293}$ and its sample/value reduction is lossless.

The stronger $2^{-47}$ public announcement \cite{announcement} must be
resolved before any best-known-bound claim. A separate $1/2$ result
\cite{abhm} uses an order fixed independently of samples and values;
it does not supply the full-seed quantifier above. The draft and the
source dependency audit received a second pass by the same agent,
not independent mathematical review. No Lean build certifies this draft.

\section{A fully specified fixed-vector rule}

First fix $w\in[0,\infty)^E$ independently of the seed. A rule reveals
and sacrifices a seed-selected mask, then receives all labels in an
arbitrary order. Its target is $\E_R\min_\pi w(A(w,R,\pi))$.
Choose
\begin{equation}\label{eq:params}
 B=2^{16},\quad \kappa=2^{18},\quad t=2^{-29},\quad
 p=1-2^{-35},\quad \beta=1-B^{-1}.
\end{equation}
Choose $J$ uniformly in $\{0,\ldots,15\}$, and round positive coordinates by
\[
 u_J(x)=2^J B^{\lfloor\log_B(x/2^J)\rfloor},\qquad u_J(0)=0.
\]
Use the same $J$ for every coordinate. The branch is main with
probability $p$ and maximum otherwise. All masks, parity, branch, and
offset coins are independent and drawn initially.

The maximum branch reveals a fair mask and accepts the first unmasked
positive nonloop with priority exceeding every revealed positive
nonloop, then stops. Higher rounded weight means higher priority,
with the fixed label order resolving ties. An empty revealed threshold
is below every positive nonloop. If $m_*$ is the greatest rounded
nonloop weight, this branch has expected worst-order reward at least
$m_*/4$: it succeeds whenever the highest-priority nonloop is unmasked
and the second is masked. With just one positive nonloop it succeeds
with probability $1/2$.

In the main branch draw masks $H,D,C,T$ with respective rates
$1/2,1/4,t,t$, independently across labels. Reveal and sacrifice
$H\cup D\cup C$. Define the analytical candidate set
\[
 U=\{e\notin H:u_e>0,\ e\notin\cl(H_{>e})\},
 \qquad Y=I^*\setminus H,
\]
where $I^*$ is priority greedy on the whole rounded vector and $H_{>e}$
contains the higher-priority labels in $H$. The candidate test uses only
the revealed $H$ values and the value at the tested label. The source's
candidate-mass lemma gives
\begin{equation}\label{eq:filter}
 Y\subseteq U,\qquad \E_H u(Y)=W/2,\qquad \E_H u(U)\le W,
 \qquad W=\OPT_M(u).
\end{equation}
Indeed the probability contribution of a candidate equals that of a
label chosen by greedy on $H$, conditional on its predecessors in $H$.

For a fixed offset divide weights by $2^J$ in the analysis, so levels
are $B^i$. Put $U_i=\{e\in U:u_e=B^i\}$, $n_i=|U_i|$,
$Y_i=Y\cap U_i$. List only levels intersecting $D$, in increasing weight
order, as $G_1,\ldots,G_m$. Write $D_h=D\cap G_h$ and likewise for
$C_h,T_h$. Their revealed parts $D_h,C_h$ suffice for preparation;
the full groups are analysis objects.

For a flat $P$, let $Q_h(P)$ be the largest maximizing flat of
$|D_h\cap Z|-\kappa\rk(Z)$ over flats $Z\supseteq P$.
For nonflat arguments first take closure. Set $a_h=-3h$ and let
$Q_{h,k}=Q_h$ for $k\ge a_h$, the identity otherwise. Starting from
$F_0(k)=L$ for $k<0$ and $F_0(k)=E$ otherwise, set
\begin{equation}\label{eq:paths}
 X_h(k)=Q_{h,k}(F_{h-1}(k)),\qquad
 F_h(k)=\cl\bigl(X_h(k)\cup(C_h\cap X_h(k+1))\bigr).
\end{equation}
These paths are computed in increasing $h$. All relevant times lie in
a finite interval of length $O(n)$. Let $\widehat F=F_m$; reject
everything if $m=0$. Draw a fair parity $\xi\in\{0,1\}$.

For an arriving label $e$, reject sacrificed labels, failed candidates,
and unlisted levels. For its listed index $h$, find its nominal birth
$b_h(e)=\min\{k:e\in X_h(k)\}$. It is eligible precisely when
\[
 e\in T,\quad b_h(e)\ge a_h+2,\quad b_h(e)\equiv\xi\pmod2,
 \quad e\notin\widehat F(b_h(e)-2).
\]
At each endpoint $b$ maintain a set $A_b$. Accept an eligible $e$ of
birth $b$ exactly when
$\rk(\{e\}\mid\widehat F(b-2)\cup A_b)=1$.

The source density and path lemmas hold for every $\kappa>0$:
$Q_h$ is increasing and extensive, and
\begin{align}
 |D_h\cap(\cl(Q_h(P)\cup Z)\setminus Q_h(P))|
   &\le\kappa\rk(Z\mid Q_h(P)),\label{eq:density}\\
 F_{h-1}(k)\subseteq X_h(k)\subseteq F_h(k)&\subseteq X_h(k+1).
 \label{eq:nesting}
\end{align}
They follow from supermodularity of the density objective on flats
and from the path recursion. Nonloops are born at $a_h$ or in
$[a_h+2,0]$. Each path has generators $J_h\subseteq D_h$ with
increasing subsets $J_h(k)$ such that
\begin{equation}\label{eq:generators}
 X_h(k)=\cl(F_{h-1}(k)\cup J_h(k)),\qquad
 |J_h|\le |D_h|/\kappa.
\end{equation}
The rank costs telescope as new density generators are inserted.
These facts require no large lower bound on $\kappa$.
By \eqref{eq:nesting}, successive layer bases contain all earlier
layers, proving independence at every arrival prefix. Finite rank
operations, fixed tie rules, and countably many measurable rounding
levels prove measurability.

\section{Safety with a larger exit cutoff}

Fix $H,D$ and a listed index $h$. Define
\[
 K_{h,b}=\bigcup_{j<h}\bigl(T_j\cap(X_j(b)\setminus X_j(b-1))\bigr).
\]
It contains all lower-weight eligible competitors in layer $b$ and
is contained in $F_{h-1}(b)$. A label $d\in D_h\cap Y$ is safe if its
birth $b$ is in $[a_h+2,0]$, has parity $\xi$, and
$d\notin Q_h(F_{h-1}(b-2)\cup K_{h,b})$. Let $\sigma_h$ count them.

The source's reverse-exposure construction queries departing guard
coordinates once each, in predictable batches. Conditional on a
transcript just before a step, the surviving coordinates retain their
product law. If a nonloop is still present, let $q_{b-1}$ be its
conditional exit probability at the next reverse step; put it zero
after exit. The disjoint-query-certificate lemma bounds the expected
number of survivals in a hazard band $[q_0,2q_0)$ by
$\log(1/q_0)$. Adding the possible exit step and summing dyadic bands
therefore gives, for an integer cutoff exponent $m_0$,
\begin{equation}\label{eq:hazard}
 \Prb(q_{\rm exit}<2^{-m_0})
 \le 2^{1-m_0}((m_0+2)\log2+1).
\end{equation}
This is the source argument with its final cutoff left symbolic. It
uses a finite query horizon and an increasing Boolean event; it does
not use the original numerical $t,B,\kappa$.

The positive series
$\log2=2\sum_{j\ge0}((2j+1)3^{2j+1})^{-1}$ implies
$2/3<\log2\le25/36<7/10$. For the upper bound, keep the first term
and bound subsequent denominators $2j+1$ by $3$.
Taking $m_0=4$ in \eqref{eq:hazard} gives a failure probability at most
$13/20$. Every nonloop exits, so
\begin{equation}\label{eq:square}
 \sum_b\E[\one_{\{d\in X_h(b)\}}q_{b-1}^2\mid H,D]
 =\E[q_{\rm exit}\mid H,D]\ge\frac7{320}.
\end{equation}

The source hybrid transition replaces just the departing guard bits
by their independent test bits. Conditional on the complete ordinary
transcript, unqueried guard bits and these test bits jointly have the
original product law. Since guard and test rates agree, the hybrid
exit has conditional probability $q_{b-1}$, and the probability of
ordinary and hybrid exit together is
$\one_{\{d\in X_h(b)\}}q_{b-1}^2$ given the pre-step transcript.
The hybrid output contains
$Q_h(F_{h-1}(b-2)\cup K_{h,b})$ for valid $b$; hence double exit
implies safety before parity. This is an endpointwise argument,
not an assertion that different endpoint experiments are independent.

The only excluded birth is $a_h$; none occurs at $a_h+1$.
Since $Y$ is independent, at most
$\rk(Q_h(L))\le |D_h|/\kappa$ of its sampled labels have this birth.
Subtract this loss from \eqref{eq:square} and then average fair parity:
\[
 \E[\sigma_h\mid H,D]
 \ge\frac7{640}|D_h\cap Y|-\frac{|D_h|}{2\kappa}.
\]
Extend $\sigma_i$ by zero on unlisted true groups. Averaging $D$ without
conditioning on listing gives
\begin{equation}\label{eq:sigma}
 \E[\sigma_i\mid H]\ge\gamma|Y_i|-\frac{n_i}{8\kappa},
 \qquad\gamma=\frac7{2560}.
\end{equation}

\section{Residual entropy and nominal rank}

Let $\mathcal B_h=\{b:a_h+2\le b\le0,\ b\equiv\xi\pmod2\}$.
For $b\in\mathcal B_h$ put
\[
 P_{h,b}=G_h\cap(X_h(b)\setminus X_h(b-1)),\quad O=D\cup C,
 \quad S_{h,b}=P_{h,b}\setminus O,
\]
and set $P_{h,b}=\varnothing$ for other endpoints of this parity.
The nominal rank is
\[
 z_h=\sum_b\rk(S_{h,b}\mid X_h(b-2)\cup K_{h,b}).
\]
Choose a label-ordered basis extension $Z_b\subseteq S_{h,b}$ for
each summand, so $\sum_b|Z_b|=z_h$. The residual
\[
 \mathcal R_h=\bigcup_b\left(P_{h,b}\setminus
  \cl(X_h(b-2)\cup K_{h,b}\cup Z_b)\right)
\]
is contained in $O$. Applying \eqref{eq:density} with the larger base
$Q_h(F_{h-1}(b-2)\cup K_{h,b})$ gives the deterministic inequality
\begin{equation}\label{eq:residual}
 \sigma_h\le\kappa z_h+|\mathcal R_h|.
\end{equation}

We use the following precise interface of the source marked-pivot
and residual-family lemmas. Conditional on $H$, parity, and all
$D,C,T$ bits outside a true group $U_i$ of size $n\ge\kappa$, there
is a fixed family $\mathfrak R_i$ containing every residual on every
listed outcome with $z_h\le n/\kappa$. For
$q=\lfloor n/\kappa\rfloor$ and
$S_d(n)=\sum_{\ell=0}^d\binom n\ell$, it satisfies
\begin{equation}\label{eq:family}
 |\mathfrak R_i|\le S_q(n)^2\bigl(4^{2q}S_{2q}(n)\bigr)^2.
\end{equation}
No focal bit or listing event is in this conditioning. The family
allows at most $q$ density generators and $q$ witness labels; two
marked-pivot sequences encode valid births and failure of spanning.
The source marked-pivot lemma bounds $s$ movable occurrences by
$4^sS_s(n)$ regardless of the timeline length. That combinatorial
lemma, including its matroid and hyperplane-arrangement proof, remains
an external dependency of this note.

For $0<\alpha\le1/2$,
$\alpha^{\alpha n}S_{\lfloor\alpha n\rfloor}(n)
\le(1+\alpha)^n\le e^{\alpha n}$.
Use $\alpha=1/\kappa,2/\kappa$ in \eqref{eq:family}, keeping the
floors, to obtain
\[
 \log|\mathfrak R_i|\le\frac n\kappa
       (6\log\kappa+6+4\log2).
\]
This sharper expression is already inside the source proof, before
its relaxation to $1000\log\kappa$. For $\kappa=2^{18}$ and
$\log2>2/3$, it implies
$|\mathfrak R_i|\le2^{121n/\kappa}$.

Under precisely the conditioning above, each focal coordinate belongs
to $O$ independently with probability $1/4+3t/4\le1/2$.
Set $\delta=139/\kappa$. A union bound gives
\[
 \Prb(\exists R_0\in\mathfrak R_i:
 R_0\subseteq O,\ |R_0|>\delta n)
 \le2^{-18n/\kappa}\le2^{-18}.
\]
Set $e_i=|\mathcal R_h|$ on listed outcomes with $z_i\le n/\kappa$,
and zero otherwise. Then
$\E[e_i\mid H]\le(\delta+2^{-18})n=140n/\kappa$.
Extend $z_i$ by zero on unlisted groups. The inequality
$\sigma_i\le\kappa z_i+e_i$ holds also on unlisted outcomes and
when $z_i>n/\kappa$, since $\sigma_i\le n$.
Together with \eqref{eq:sigma}, this proves, for $n_i\ge\kappa$,
\begin{equation}\label{eq:z}
 \E[z_i\mid H]\ge\frac{\gamma|Y_i|-\epsilon n_i}{\kappa},
 \qquad\epsilon=\frac{1121}{2^{21}}.
\end{equation}

\section{Order-uniform accounting and branch mixing}

Define
\[
 \lambda_h=\sum_b\rk(T_h\cap S_{h,b}\mid\widehat F(b-2)\cup K_{h,b}).
\]
With all coins fixed, the source order-uniform rank lemma gives
$N_{\ge h}(\pi)\ge\lambda_h$ for every permutation, where $N_{\ge h}$
counts accepts of weight at least that of $G_h$. Indeed, the final
greedy set of a layer spans its eligible labels; its selected lower
weight labels lie in $K_{h,b}$. Contracting them proves the bound.

One common set of generators for the later paths is
$J'=C_h\cup\bigcup_{j>h}(C_j\cup J_j)$, of size at most
\[
 c_h=|C_h|+\sum_{j>h}(|C_j|+|D_j|/\kappa).
\]
It satisfies $X_h(k)\subseteq\widehat F(k)
\subseteq\cl(X_h(k)\cup J')$ simultaneously for all $k$.
The resulting rank losses at chosen-parity endpoints telescope, giving
\[
 \sum_b\rk(S_{h,b}\mid\widehat F(b-2)\cup K_{h,b})\ge z_h-c_h.
\]
Specifically each loss is at most
$\rk(J'\mid X_h(b-2))-\rk(J'\mid X_h(b))$, so the same generators
are charged only once.

Condition on $H,D,C,\xi$ and all $T$ bits outside $G_h$.
The summand sets and bases are fixed and focal $T$ remains an independent
rate-$t$ mask. Intersecting a fixed basis extension with $T_h$ shows
$\E[\lambda_h\mid\cdot]\ge t(z_h-c_h)$.
Extend $\lambda_i,c_i$ by zero on unlisted true groups. Bounding costs
by counts on all true groups before averaging gives
\[
 \E[c_i\mid H]\le t n_i+(t+1/(4\kappa))\sum_{\ell>i}n_\ell.
\]
Thus \eqref{eq:z} implies
\begin{equation}\label{eq:lambda}
 \E[\lambda_i\mid H]\ge\frac t\kappa
 \left[\gamma|Y_i|-(\epsilon+t\kappa)n_i
              -(t\kappa+1/4)\sum_{\ell>i}n_\ell\right]
 \quad(n_i\ge\kappa).
\end{equation}
No chosen arrival order has been conditioned upon.

Each accepted weight $B^\ell$ is counted in cumulative level counts
by at most $\sum_{i\le\ell}B^i=B^\ell/\beta$. Consequently
\begin{equation}\label{eq:min}
 \min_\pi u(A_{\rm main}(\pi))\ge
 \beta\sum_{i:n_i\ge\kappa}B^i\lambda_i.
\end{equation}
Also
\[
 \sum_iB^i\sum_{\ell>i}n_\ell\le\frac{u(U)}{B-1},\qquad
 \sum_{i:n_i<\kappa}B^i|Y_i|\le\frac{\kappa m_*}{\beta}.
\]
Multiply \eqref{eq:lambda} by levels, use \eqref{eq:min}, and average
\eqref{eq:filter}. Restoring the factor $2^J$ when needed gives
\begin{equation}\label{eq:affine}
 \E\min_\pi u(A_{\rm main}(\pi))\ge AW-t\gamma m_*,\qquad
 A=\frac{\beta t}{\kappa}s,
\end{equation}
where
\begin{equation}\label{eq:slack}
 s=\frac\gamma2-\epsilon-t\kappa-
          \frac{t\kappa+1/4}{B-1}
   =\frac{3120271}{9162457088}>\frac1{3072}.
\end{equation}
The affine bound may be negative for a particular vector; that causes
no problem. Mix in the maximum branch:
\[
 \E\min_\pi u(A(\pi))\ge
 pAW+\left(\frac{1-p}{4}-pt\gamma\right)m_*\ge pAW,
\]
because $(1-p)/4=t/256$ and $\gamma=7/2560<1/256$.
Both minima occur before averaging the respective seeds, so the
adversary may select its order after seeing the branch.

\section{Finite offsets and lossless one-sample transfer}

For any $x>0$, let $k=\lfloor\log_2x\rfloor$. As $J$ ranges over
the sixteen phases, rounded exponents are exactly $k,k-1,\ldots,k-15$.
Therefore
\[
 \E_J u_J(x)=\frac{2^{k+1}\beta}{16}\ge\frac\beta{16}x.
\]
The same inequality holds at zero. Fix an optimum independent set
for $w$ before drawing $J$, and sum to get
$\E_J\OPT_M(u_J)\ge\beta\OPT_M(w)/16$.
Conditional on any phase the remaining seed has its original product
law. True accepted weights dominate rounded weights, so
\begin{equation}\label{eq:coefficient}
 \E_{J,R}\min_\pi w(A(J,R,\pi))
 \ge C\OPT_M(w),\qquad C=\frac{p\beta^2ts}{16\kappa}>2^{-63}.
\end{equation}
Indeed $p\beta^2>3/4$, $s>1/3072$, and $t/\kappa=2^{-47}$,
whose product in \eqref{eq:coefficient} is $2^{-63}$.
The exact coefficient is
\[
 C=\frac{105391776668216973667425}
 {696898287454081973172991196020261297061888}
 \approx1.39485\cdot2^{-63}.
\]
Visibility of the phase does not change this argument: the minimum
over orders was already taken at each complete seed.

For the prophet rule, let $B_0(R)$ be its sacrificed mask and define
$w_e^{\rm glue}=S_e$ on $B_0(R)$ and $V_e$ elsewhere.
Conditional on the complete seed, this vector has the product law
of $V$, independent of that seed. Run the hidden-vector rule using
samples only on $B_0$, reject those labels, and use arriving values
on its complement. For each $S,V,R$ its reward under any adversarial
order is at least the minimum over all permutations on the glued
vector. Integrating \eqref{eq:coefficient} proves the proposed theorem.
Unused samples and values may influence the adversary without affecting
this lower bound. Feasibility bounds actual reward by $\OPT_M(V)$;
nonnegative integration needs no extra moment hypothesis on loop values.

Predraw four phase bits, 35 branch bits (choose maximum only when
all are zero), one parity bit, and per label one $H$, two $D$, 29 $C$,
29 $T$, and one maximum-mask bit. This totals $62n+40$ fair bits.
The finite rule can precompute rank and closure tables over all subsets
and enumerate flats for each density map. That implementation may use
exponential work, queries, and storage. Arbitrary real observations
remain exact in the source measurable model; rational implementations
pay input-bit costs for exact level comparisons. Output work is at
least the number of accepted labels. No efficiency or precision
improvement is asserted. Empty ground sets, rank zero, loops, ties,
zero values, and arbitrary integer weight exponents are included.

\section{Evidence and release status}

The accompanying standard-library Python certificate recomputes every
scalar displayed above with exact rational arithmetic. Separate checks
enumerate all labeled matroids on up to four elements for candidate,
density, maximum-branch, and rounding facts. An online-rule checker
enumerates small complete seed states, takes the worst arrival order
before averaging, and checks prefix feasibility, path generators,
conditional safety, later-path rank cost, and order-uniform rank bounds.
Small density thresholds and a larger test rate exercise nontrivial
transitions; those tests do not measure performance at the proposed
parameters or prove a universal matroid theorem.

The quantitative derivation has been audited against the source by
the same agent that wrote it. Remaining review targets include the
parameter-uniform reverse/hybrid argument and the source residual
encoding under its precise conditioning. A separate analytic note shows
that further searching our conservative scalar-certificate family
cannot reach the announced $2^{-47}$; this is a limitation of that
certificate, not an impossibility theorem for the algorithm or problem.
This version is preparatory work in the current repository. Authorship,
independent review, novelty comparison, and a public release decision
remain to be settled before a standalone announcement.

\begin{thebibliography}{9}
\bibitem{source} OpenAI.
\emph{One Sample Suffices for Matroid Prophet Inequalities against an
Almighty Adversary}. September 23, 2026 manuscript, retrieved at
commit \texttt{fd4aeeb2ee4fc729c18d98444fed42fd0529eeeb}.
\href{https://github.com/openai/math/blob/fd4aeeb2ee4fc729c18d98444fed42fd0529eeeb/preprints/One-Sample-Suffices-for-Matroid-Prophet-Inequalities-against-an-Almighty-Adversary-September-23-2026/final.pdf}{Pinned manuscript}.
Relevant source lemma names: Candidate mass; Density expansion; Path
properties; A generator budget for the entire path; Disjoint-query
certificates; Exit mass; Safe sampled labels; Marked pivots; A uniform
family of residuals; Order-uniform rank bound; Telescoping rank cost;
Independent test thinning; One-sample simulation.
\bibitem{announcement} Daniel (@firsttimelifer).
\href{https://x.com/firsttimelifer/status/2108282829809541468}{Public
announcement of a $2^{-47}$ guarantee}, October 8, 2026.
Proof and exact assumptions unavailable in this audit.
\bibitem{abhm} Abdi, Banihashem, Hajiaghayi, and Mittal.
\emph{On the Strong Matroid Secretary Conjecture and Beyond}.
\href{https://arxiv.org/abs/2609.19118v1}{arXiv:2609.19118v1},
September 16, 2026. Theorem 1.3 and its fixed-order model.
\end{thebibliography}
\end{document}
